\newpage

\section{Selected Technology}

\subsection{Technology Stack}

The proposed e-commerce platform utilizes a modern, scalable technology stack designed to support concurrent users, complex transactions, and real-time data synchronization:

\subsubsection{Database Management System}

\textbf{MySQL Server 8.0} is selected as the primary DBMS due to its:
\begin{itemize}
    \item \textbf{ACID Compliance}: Ensures data integrity across multi-seller transactions
    \item \textbf{Foreign Key Constraints}: Enforces referential integrity across 14 entities
    \item \textbf{Triggers and Stored Procedures}: Enables business logic automation (auto-generated codes, inventory management, rating calculations)
    \item \textbf{Recursive Queries}: Supports hierarchical category traversal
    \item \textbf{JSON Support}: Optional future extensibility for product metadata
\end{itemize}

\subsubsection{Backend Framework}

\textbf{Node.js with Express.js} provides:
\begin{itemize}
    \item Event-driven, non-blocking I/O for handling multiple concurrent requests
    \item RESTful API design for clean separation between frontend and backend
    \item Middleware support for authentication, validation, and error handling
    \item Connection pooling for efficient database access
\end{itemize}

\subsubsection{Frontend Framework}

\textbf{React with Vite} offers:
\begin{itemize}
    \item Component-based architecture for reusable UI elements
    \item State management via Redux for complex application state
    \item Fast hot module replacement during development
    \item Optimized production builds with tree-shaking
\end{itemize}

\subsubsection{Authentication}

\textbf{JWT (JSON Web Tokens)} for stateless authentication:
\begin{itemize}
    \item Secure token-based sessions
    \item Role-based access control (Customer, Shop, Admin)
    \item Token expiration and refresh mechanisms
\end{itemize}

\section{Database Implementation}

\subsection{Account Management}

\subsubsection{ACCOUNT Table}

The \textbf{ACCOUNT} table serves as the central authentication and identity store for all system users. It implements table-level inheritance, with specialized entities (CUSTOMER, SHOP, ADMIN) storing additional role-specific information.

\begin{lstlisting}[language=SQL, caption=ACCOUNT Table Creation, label=lst:account]
CREATE TABLE Account (
    account_id INT AUTO_INCREMENT PRIMARY KEY,
    email VARCHAR(255) UNIQUE NOT NULL,
    password VARCHAR(255) NOT NULL,
    role ENUM('Customer','Shop','Admin') NOT NULL,
    full_name VARCHAR(255) NOT NULL,
    phone VARCHAR(20) NOT NULL CHECK (LENGTH(phone) = 10 AND phone REGEXP '^[0-9]{10}$'),
    status ENUM('Active','Ban') DEFAULT 'Active',
    created_at TIMESTAMP DEFAULT CURRENT_TIMESTAMP,
    CONSTRAINT CK_AccountRole CHECK (role IN ('Customer','Shop','Admin'))
);
\end{lstlisting}

\paragraph{General Description:}
The ACCOUNT table centralizes user authentication across all three user types. The \textbf{email} field is UNIQUE to prevent duplicate account registrations. The \textbf{phone} CHECK constraint enforces Vietnamese phone format validation at the database level. The \textbf{role} ENUM field determines user permissions (Customer can purchase, Shop can sell, Admin supervises). Status field enables account suspension without deletion.

\subsubsection{CUSTOMER Table}

\begin{lstlisting}[language=SQL, caption=CUSTOMER Table Creation, label=lst:customer]
CREATE TABLE Customer (
    customer_id INT PRIMARY KEY,
    customer_code VARCHAR(20) UNIQUE,
    address VARCHAR(255),
    add_phone VARCHAR(20) CHECK (LENGTH(add_phone) = 10 AND add_phone REGEXP '^[0-9]{10}$'),
    total_spent DECIMAL(15,2) DEFAULT 0,
    total_order INT DEFAULT 0,
    FOREIGN KEY (customer_id) REFERENCES Account(account_id) ON DELETE CASCADE
);
\end{lstlisting}

\paragraph{General Description:}
The CUSTOMER table extends ACCOUNT with purchase-specific attributes. The \textbf{customer\_code} is auto-generated via trigger with 'CUST' prefix, enabling easy customer identification in reports. \textbf{Address} is mandatory for order fulfillment. \textbf{TotalSpent} and \textbf{TotalOrder} are derived attributes maintained by triggers on order completion, supporting customer loyalty analysis. The UNIQUE constraint on customer\_code and index on total\_spent enable efficient customer ranking and loyalty program implementation.

\subsubsection{SHOP Table}

\begin{lstlisting}[language=SQL, caption=SHOP Table Creation, label=lst:shop]
CREATE TABLE Shop (
    shop_id INT PRIMARY KEY,
    shop_code VARCHAR(20) UNIQUE,
    shop_name VARCHAR(255) NOT NULL,
    shop_phone VARCHAR(20) CHECK (LENGTH(shop_phone) = 10 AND shop_phone REGEXP '^[0-9]{10}$'),
    address_shop VARCHAR(255),
    rating DECIMAL(3,2) DEFAULT 0 CHECK (rating BETWEEN 0 AND 5),
    shop_status ENUM('Open','Temporarily Close','Closed') DEFAULT 'Open',
    FOREIGN KEY (shop_id) REFERENCES Account(account_id) ON DELETE CASCADE
);
\end{lstlisting}

\paragraph{General Description:}
The SHOP table models seller entities with business-specific attributes. \textbf{shop\_code} is auto-generated with 'SHOP' prefix for seller identification. \textbf{shop\_status} controls whether a seller can list products (Open) or is inactive. \textbf{rating} is auto-calculated from customer reviews via trigger, providing trust metrics. Indexes on shop\_code and rating enable efficient shop searching and ranking by customer ratings.

\subsubsection{ADMIN Table}

\begin{lstlisting}[language=SQL, caption=ADMIN Table Creation, label=lst:admin]
CREATE TABLE Admin (
    admin_id INT PRIMARY KEY,
    role ENUM('Core','Moderator','Support') NOT NULL,
    note TEXT,
    FOREIGN KEY (admin_id) REFERENCES Account(account_id) ON DELETE CASCADE
);
\end{lstlisting}

\paragraph{General Description:}
The ADMIN table represents platform supervisors with three access levels. \textbf{Core} admins have full system access, \textbf{Moderator} admins handle dispute resolution and content moderation, while \textbf{Support} admins assist customers. The note field allows tracking of admin activities and responsibilities. Unlike CUSTOMER and SHOP, ADMIN entities do not engage in transactions.

\subsection{Product Catalog}

\subsubsection{CATEGORY Table}

\begin{lstlisting}[language=SQL, caption=CATEGORY Table Creation, label=lst:category]
CREATE TABLE Category (
    category_id INT AUTO_INCREMENT PRIMARY KEY,
    category_name VARCHAR(100) NOT NULL,
    parent_category_id INT,
    FOREIGN KEY (parent_category_id) REFERENCES Category(category_id) ON DELETE SET NULL
);
\end{lstlisting}

\paragraph{General Description:}
The CATEGORY table implements a self-referential recursive relationship enabling hierarchical product organization (e.g., Electronics $\rightarrow$ Phones $\rightarrow$ Smartphones). The \textbf{parent\_category\_id} can be NULL for root-level categories. The recursive FK with ON DELETE SET NULL ensures child categories remain accessible if a parent is deleted..

\subsubsection{PRODUCT Table}

\begin{lstlisting}[language=SQL, caption=PRODUCT Table Creation, label=lst:product]
CREATE TABLE Product (
    product_id INT AUTO_INCREMENT PRIMARY KEY,
    shop_id INT NOT NULL,
    category_id INT NOT NULL,
    product_name VARCHAR(255) NOT NULL,
    description TEXT,
    image VARCHAR(255),
    rating DECIMAL(3,2) DEFAULT 0,
    status ENUM('In stock','Out of stock') DEFAULT 'In stock',
    created_at TIMESTAMP DEFAULT CURRENT_TIMESTAMP,
    FOREIGN KEY (shop_id) REFERENCES Shop(shop_id) ON DELETE CASCADE,
    FOREIGN KEY (category_id) REFERENCES Category(category_id)
);
\end{lstlisting}

\paragraph{General Description:}
The PRODUCT table represents individual product listings managed by shops. Each product belongs to exactly one SHOP (N:1 relationship with ON DELETE CASCADE to remove products when shop is deleted) and one CATEGORY (N:1 with ON DELETE RESTRICT to preserve category structure). The \textbf{rating} is auto-calculated from customer reviews via trigger. \textbf{status} tracks inventory availability and is auto-updated by triggers when stock reaches zero. The updated\_at timestamp enables tracking product modification history.

\subsubsection{PRODUCTITEM Table}

\begin{lstlisting}[language=SQL, caption=PRODUCTITEM Table Creation, label=lst:productitem]
CREATE TABLE ProductItem (
    item_id INT AUTO_INCREMENT PRIMARY KEY,
    product_id INT NOT NULL,
    shop_id INT NOT NULL,
    color VARCHAR(50) NOT NULL,
    type VARCHAR(50) NOT NULL,
    price DECIMAL(15,2) NOT NULL CHECK(price >= 0),
    stock INT NOT NULL CHECK(stock >= 0),
    image_url VARCHAR(255),
    FOREIGN KEY (product_id) REFERENCES Product(product_id) ON DELETE CASCADE,
    FOREIGN KEY (shop_id) REFERENCES Shop(shop_id) ON DELETE CASCADE
);
\end{lstlisting}

\paragraph{General Description:}
The PRODUCTITEM (Variant) table represents purchasable variations of products. The UNIQUE constraint on (product\_id, color, type) prevents duplicate variants of the same product. The \textbf{stock} field tracks inventory levels and is decremented by triggers on order placement, incremented on order cancellation. The \textbf{price} field is stored for each variant, enabling different pricing for sizes/colors.

\subsection{Shopping Cart}

\subsubsection{CART Table}
\begin{lstlisting}[language=SQL, caption=CART Table Creation, label=lst:cart]
CREATE TABLE Cart (
    cart_id INT AUTO_INCREMENT PRIMARY KEY,
    customer_id INT UNIQUE NOT NULL,
    created_at TIMESTAMP DEFAULT CURRENT_TIMESTAMP,
    updated_at TIMESTAMP DEFAULT CURRENT_TIMESTAMP ON UPDATE CURRENT_TIMESTAMP,
    FOREIGN KEY (customer_id) REFERENCES Customer(customer_id) ON DELETE CASCADE
);
\end{lstlisting}

\paragraph{General Description:}
The CART table implements a 1:1 relationship with CUSTOMER, ensuring each customer has exactly one shopping cart. The UNIQUE constraint on customer\_id enforces this one-to-one relationship. The updated\_at timestamp tracks cart modification time, useful for session management. ON DELETE CASCADE ensures cart deletion when customer account is removed.

\subsubsection{CARTITEM Table}

\begin{lstlisting}[language=SQL, caption=CARTITEM Table Creation, label=lst:cartitem]
CREATE TABLE CartItem (
    cart_item_id INT AUTO_INCREMENT PRIMARY KEY,
    cart_id INT NOT NULL,
    item_id INT NOT NULL,
    quantity INT DEFAULT 1 CHECK(quantity >= 1),
    added_at TIMESTAMP DEFAULT CURRENT_TIMESTAMP,
    FOREIGN KEY (cart_id) REFERENCES Cart(cart_id) ON DELETE CASCADE,
    FOREIGN KEY (item_id) REFERENCES ProductItem(item_id) ON DELETE CASCADE,
    UNIQUE KEY unique_cart_item (cart_id, item_id)
);
\end{lstlisting}

\paragraph{General Description:}
The CARTITEM table implements the N:M relationship between CART and PRODUCTITEM. The UNIQUE constraint on (cart\_id, item\_id) prevents duplicate items in the same cart. The \textbf{quantity} field tracks the number of units selected, with CHECK constraint ensuring at least 1 unit. ON DELETE CASCADE ensures cart items are removed when either cart or product is deleted.


\subsection{Orders and Items}

\subsubsection{ORDER Table}

\begin{lstlisting}[language=SQL, caption=ORDER Table Creation, label=lst:order]
CREATE TABLE `Order` (
    order_id INT AUTO_INCREMENT PRIMARY KEY,
    customer_id INT NOT NULL,
    shipping_id INT NOT NULL,
    voucher_id INT,
    status ENUM('Processing','Shipped','Delivered','Cancelled') DEFAULT 'Processing',
    order_date TIMESTAMP DEFAULT CURRENT_TIMESTAMP,
    delivered_date TIMESTAMP NULL,
    shipping_address VARCHAR(255) NOT NULL,
    total_amount DECIMAL(15,2) DEFAULT 0,
    payment_method ENUM('COD','Banking','Momo','ZaloPay') NOT NULL DEFAULT 'COD',
    note TEXT,
    created_at TIMESTAMP DEFAULT CURRENT_TIMESTAMP,
    updated_at TIMESTAMP DEFAULT CURRENT_TIMESTAMP ON UPDATE CURRENT_TIMESTAMP,
    FOREIGN KEY (customer_id) REFERENCES Customer(customer_id),
    FOREIGN KEY (shipping_id) REFERENCES Shipping(shipping_id),
    FOREIGN KEY (voucher_id) REFERENCES Voucher(voucher_id) ON DELETE SET NULL,
    CONSTRAINT CK_OrderDateRange CHECK (delivered_date IS NULL OR delivered_date >= order_date)
);
\end{lstlisting}

\paragraph{General Description:}
The ORDER table records all purchase transactions. The \textbf{status} ENUM tracks order lifecycle (Processing → Shipped → Delivered or Cancelled). The \textbf{voucher\_id} is nullable (N:0..1 relationship) since orders may not use vouchers. The CHECK constraint \texttt{delivered\_date >= order\_date} enforces chronological consistency. Indexes on customer\_id, status, and order\_date enable efficient order retrieval and reporting. ON DELETE RESTRICT on shipping\_id prevents accidental deletion of shipping methods in use.

\subsubsection{ORDERITEM Table}

\begin{lstlisting}[language=SQL, caption=ORDERITEM Table Creation, label=lst:orderitem]
CREATE TABLE OrderItem (
    order_item_id INT AUTO_INCREMENT PRIMARY KEY,
    order_id INT NOT NULL,
    item_id INT NOT NULL,
    shop_id INT NOT NULL,
    quantity INT NOT NULL CHECK(quantity >= 1),
    price_at_purchase DECIMAL(15,2) NOT NULL CHECK(price_at_purchase >= 0),
    FOREIGN KEY (order_id) REFERENCES `Order`(order_id) ON DELETE CASCADE,
    FOREIGN KEY (item_id) REFERENCES ProductItem(item_id),
    FOREIGN KEY (shop_id) REFERENCES Shop(shop_id)
);
\end{lstlisting}

\paragraph{General Description:}
The ORDERITEM table represents individual line items within an order, supporting multi-seller purchases. The \textbf{price\_at\_purchase} field captures the price at transaction time, preserving historical accuracy even if product prices change. The \textbf{shop\_id} field identifies the seller responsible for fulfilling each line item. The UNIQUE constraint on (order\_id, item\_id) prevents duplicate line items. ON DELETE RESTRICT on item\_id and shop\_id ensures order history remains intact.

\subsection{Promotions}

\subsubsection{VOUCHER Table}

\begin{lstlisting}[language=SQL, caption=VOUCHER Table Creation, label=lst:voucher]
CREATE TABLE Voucher (
    voucher_id INT AUTO_INCREMENT PRIMARY KEY,
    code VARCHAR(50) UNIQUE NOT NULL,
    discount_type ENUM('Percentage','Amount') NOT NULL,
    discount_value DECIMAL(10,2) NOT NULL CHECK(discount_value >= 0),
    min_order_value DECIMAL(15,2) DEFAULT 0 CHECK(min_order_value >= 0),
    start_date DATE NOT NULL,
    expired_date DATE NOT NULL,
    usage_limit INT DEFAULT 1 CHECK(usage_limit >= 1),
    used_count INT DEFAULT 0,
    status ENUM('Active','Expired') DEFAULT 'Active',
    created_at TIMESTAMP DEFAULT CURRENT_TIMESTAMP,
    CONSTRAINT CK_VoucherDateRange CHECK (expired_date >= start_date)
);
\end{lstlisting}

\paragraph{General Description:}
The VOUCHER table manages promotional campaigns with flexible discount mechanisms. The \textbf{discount\_type} ENUM supports both percentage (e.g., 10\%) and fixed amount (e.g., 50,000 VND) discounts. The \textbf{min\_order\_value} enforces minimum purchase requirements. The \textbf{usage\_limit} and \textbf{used\_count} fields track voucher usage, preventing over-utilization. The CHECK constraint ensures used\_count never exceeds usage\_limit.

\subsection{Shipping}
\subsubsection{SHIPPING Table}

\begin{lstlisting}[language=SQL, caption=SHIPPING Table Creation, label=lst:shipping]
CREATE TABLE Shipping (
    shipping_id INT AUTO_INCREMENT PRIMARY KEY,
    name VARCHAR(100) NOT NULL,
    estimated_days INT NOT NULL CHECK(estimated_days >= 0),
    fee DECIMAL(10,2) NOT NULL CHECK(fee >= 0),
    status ENUM('Active','Inactive') DEFAULT 'Active'
);
\end{lstlisting}

\paragraph{General Description:}
The SHIPPING table defines available delivery options for orders. Each record represents a logistics partner or delivery method. The \textbf{estimated\_days} field provides expected delivery timeframes. The \textbf{fee} field enables dynamic shipping cost calculation. The \textbf{status} field controls availability (Active methods appear in checkout, Inactive ones are hidden).

\subsection{Reviews and Feedback}
\subsubsection{REVIEW Table}

\begin{lstlisting}[language=SQL, caption=REVIEW Table Creation, label=lst:review]
CREATE TABLE Review (
    review_id INT AUTO_INCREMENT PRIMARY KEY,
    customer_id INT NOT NULL,
    target_type ENUM('Shop','Product') NOT NULL,
    target_id INT NOT NULL,
    rating INT NOT NULL CHECK(rating BETWEEN 1 AND 5),
    comment TEXT NOT NULL,
    image_url VARCHAR(255),
    review_date TIMESTAMP DEFAULT CURRENT_TIMESTAMP,
    FOREIGN KEY (customer_id) REFERENCES Customer(customer_id) ON DELETE CASCADE
);

\end{lstlisting}

\section{Sample Data Insertion}

This section demonstrates sample data for testing the database schema. Data is organized by entity relationships to maintain referential integrity.

\subsection{Account Data}

\begin{lstlisting}[language=SQL, caption=Sample ACCOUNT Insert, label=lst:insert_account]
-- Insert 7 accounts: 3 customers, 3 shops, 1 admin
INSERT INTO Account (email, password, role, full_name, phone, status)
VALUES
    ('c1@gmail.com', '1234', 'Customer', 'Nguyen Van A', '0901111111', 'Active'),
    ('c2@gmail.com', '1234', 'Customer', 'Tran Thi B', '0902222222', 'Active'),
    ('c3@gmail.com', '1234', 'Customer', 'Le Hoai C', '0903333333', 'Active'),
    ('s1@gmail.com', '1234', 'Shop', 'Shop ABC', '0904444444', 'Active'),
    ('s2@gmail.com', '1234', 'Shop', 'Shop XYZ', '0905555555', 'Active'),
    ('s3@gmail.com', '1234', 'Shop', 'Tech Store', '0906666666', 'Active'),
    ('admin1@gmail.com', '1234', 'Admin', 'Admin Core', '0907777777', 'Active');
\end{lstlisting}

\subsection{Customer Data}

\begin{lstlisting}[language=SQL, caption=Sample CUSTOMER Insert, label=lst:insert_customer]
-- Triggers automatically create CUSTOMER records during INSERT
-- This example shows updating customer info with addresses
UPDATE Customer SET 
    address = '123 Nguyen Hue, District 1, HCMC',
    add_phone = '0908888888'
WHERE customer_id = 1;

UPDATE Customer SET 
    address = '456 Le Loi, District 3, HCMC',
    add_phone = '0909999999'
WHERE customer_id = 2;

UPDATE Customer SET 
    address = '789 Tran Hung Dao, Hoan Kiem, Hanoi',
    add_phone = '0911111111'
WHERE customer_id = 3;
\end{lstlisting}

\subsection{Shop Data}

\begin{lstlisting}[language=SQL, caption=Sample SHOP Insert, label=lst:insert_shop]
-- Triggers automatically create SHOP records
-- Update shop information
UPDATE Shop SET 
    shop_name = 'Shop ABC - Electronics',
    shop_phone = '0904444444',
    address_shop = '12 Vo Van Kiet, District 1',
    rating = 4.50,
    shop_status = 'Open'
WHERE shop_id = 4;

UPDATE Shop SET 
    shop_name = 'Shop XYZ - Fashion',
    shop_phone = '0905555555',
    address_shop = '45 Hai Ba Trung, District 3',
    rating = 0,
    shop_status = 'Open'
WHERE shop_id = 5;

UPDATE Shop SET 
    shop_name = 'Tech Store - Premium',
    shop_phone = '0906666666',
    address_shop = '67 Nguyen Dinh Chieu, District 7',
    rating = 5.00,
    shop_status = 'Open'
WHERE shop_id = 6;
\end{lstlisting}

\subsection{Category Data}

\begin{lstlisting}[language=SQL, caption=Sample CATEGORY Insert, label=lst:insert_category]
INSERT INTO Category (category_name, parent_category_id)
VALUES
    ('Electronics', NULL),              -- Root: ID 1
    ('Fashion', NULL),                  -- Root: ID 2
    ('Phones', 1),                      -- Child of Electronics
    ('Laptops', 1),                     -- Child of Electronics
    ('Clothes', 2),                     -- Child of Fashion
    ('Shoes', 2);                       -- Child of Fashion
\end{lstlisting}

\subsection{Product and ProductItem Data}

\begin{lstlisting}[language=SQL, caption=Sample PRODUCT Insert, label=lst:insert_product]
-- Insert products from different shops
INSERT INTO Product (shop_id, category_id, product_name, 
                     description, image, status)
VALUES
    (4, 3, 'iPhone 15 Pro',
     'Latest Apple flagship with A17 Pro chip',
     'iphone15pro.jpg', 'In stock'),
    
    (4, 3, 'Samsung Galaxy S24',
     'Latest Samsung with AI features',
     'galaxyS24.jpg', 'In stock'),
    
    (5, 5, 'Premium Cotton T-Shirt',
     'High-quality 100% cotton shirt',
     'tshirt.jpg', 'In stock'),
    
    (6, 4, 'MacBook Air M3',
     'Lightweight laptop for professionals',
     'macbookm3.jpg', 'In stock'),
    
    (6, 3, 'Sony Headset XM5',
     'Noise-cancelling wireless headphones',
     'headset.jpg', 'In stock');
\end{lstlisting}

\begin{lstlisting}[language=SQL, caption=Sample PRODUCTITEM Insert, label=lst:insert_productitem]
-- Insert product variants (color, size, capacity)
INSERT INTO ProductItem (product_id, shop_id, color, type, 
                        price, stock, image_url)
VALUES
    -- iPhone 15 Pro variants (Product 1, Shop 4)
    (1, 4, 'Space Black', '128GB', 1099000, 50,
     'iphone15pro-black-128.jpg'),
    
    (1, 4, 'Blue Titanium', '256GB', 1199000, 30,
     'iphone15pro-blue-256.jpg'),
    
    -- Samsung Galaxy S24 variants (Product 2, Shop 4)
    (2, 4, 'Phantom Black', '256GB', 920000, 40,
     'galaxyS24-black-256.jpg'),
    
    (2, 4, 'Amber Gold', '512GB', 1099000, 25,
     'galaxyS24-gold-512.jpg'),
    
    -- T-Shirt variants (Product 3, Shop 5)
    (3, 5, 'White', 'Size M', 150000, 100,
     'tshirt-white-m.jpg'),
    
    (3, 5, 'Black', 'Size L', 150000, 80,
     'tshirt-black-l.jpg'),
    
    -- MacBook Air M3 variants (Product 4, Shop 6)
    (4, 6, 'Silver', '16GB RAM', 2299000, 15,
     'macbookm3-silver-16gb.jpg'),
    
    (4, 6, 'Space Gray', '24GB RAM', 2599000, 10,
     'macbookm3-gray-24gb.jpg'),
    
    -- Sony Headset variants (Product 5, Shop 6)
    (5, 6, 'Black', 'Standard', 4990000, 60,
     'sonyxm5-black.jpg'),
    
    (5, 6, 'Silver', 'Standard', 4990000, 40,
     'sonyxm5-silver.jpg');
\end{lstlisting}

\subsection{Shipping and Voucher Data}

\begin{lstlisting}[language=SQL, caption=Sample SHIPPING Insert, label=lst:insert_shipping]
INSERT INTO Shipping (name, estimated_days, fee, status)
VALUES
    ('GHN Express', 2, 15000, 'Active'),
    ('J&T Standard', 3, 12000, 'Active'),
    ('Shopee Express', 1, 20000, 'Active'),
    ('Ninja Van Economy', 4, 9000, 'Active');
\end{lstlisting}

\begin{lstlisting}[language=SQL, caption=Sample VOUCHER Insert, label=lst:insert_voucher]
INSERT INTO Voucher (code, discount_type, discount_value, 
                     min_order_value, start_date, expired_date, 
                     usage_limit, status)
VALUES
    -- 10% discount, min 500K
    ('DISCOUNT10', 'Percentage', 10, 500000,
     '2025-01-01', '2025-12-31', 1000, 'Active'),
    
    -- Fixed 100K discount, min 1M
    ('SAVE100K', 'Amount', 100000, 1000000,
     '2025-06-01', '2025-12-31', 500, 'Active'),
    
    -- New Year promotion: 5% off
    ('NEWYEAR2025', 'Percentage', 5, 200000,
     '2025-12-01', '2026-01-31', 2000, 'Active'),
    
    -- Expired voucher (for testing)
    ('OLDCODE', 'Percentage', 15, 300000,
     '2024-01-01', '2024-12-31', 100, 'Expired');
\end{lstlisting}

\subsection{Order Transaction Data}

\begin{lstlisting}[language=SQL, caption=Sample ORDER Insert, label=lst:insert_order]
-- Note: Cart and CartItem are auto-created by triggers
-- Orders created from shopping carts

-- Order 1: Customer 1, 2 items from Shop 4
INSERT INTO \`Order\` (customer_id, shipping_id, voucher_id, status,
                      shipping_address, total_amount, 
                      payment_method, note)
VALUES
    (1, 1, 1, 'Delivered',
     '123 Nguyen Hue, District 1, HCMC', 2098000,
     'Banking', 'Received in good condition');

-- Insert order items for Order 1
INSERT INTO OrderItem (order_id, item_id, shop_id, quantity, 
                       price_at_purchase)
VALUES
    (1, 1, 4, 1, 1099000),  -- iPhone 15 Pro Black 128GB
    (1, 2, 4, 1, 1199000);  -- iPhone 15 Pro Blue 256GB

-- Order 2: Customer 2, cancelled order
INSERT INTO \`Order\` (customer_id, shipping_id, voucher_id, status,
                      shipping_address, total_amount, 
                      payment_method, note)
VALUES
    (2, 2, NULL, 'Cancelled',
     '456 Le Loi, District 3, HCMC', 932000,
     'COD', 'Cancelled by customer');

INSERT INTO OrderItem (order_id, item_id, shop_id, quantity, 
                       price_at_purchase)
VALUES
    (2, 3, 4, 1, 920000);   -- Galaxy S24 Black 256GB

-- Order 3: Customer 3, delivered
INSERT INTO \`Order\` (customer_id, shipping_id, voucher_id, status,
                      shipping_address, total_amount, 
                      payment_method, note)
VALUES
    (3, 3, 2, 'Delivered',
     '789 Tran Hung Dao, Hanoi', 170000,
     'Momo', 'Please leave at reception');

INSERT INTO OrderItem (order_id, item_id, shop_id, quantity, 
                       price_at_purchase)
VALUES
    (3, 5, 5, 1, 150000);   -- T-Shirt White Size M

-- Order 4: Customer 1, shipped
INSERT INTO \`Order\` (customer_id, shipping_id, voucher_id, status,
                      shipping_address, total_amount, 
                      payment_method, note)
VALUES
    (1, 4, NULL, 'Shipped',
     '123 Nguyen Hue, District 1, HCMC', 5009000,
     'ZaloPay', NULL);

INSERT INTO OrderItem (order_id, item_id, shop_id, quantity, 
                       price_at_purchase)
VALUES
    (4, 9, 6, 1, 4990000);  -- Sony Headset Black
\end{lstlisting}

\subsection{Review Data}

\begin{lstlisting}[language=SQL, caption=Sample REVIEW Insert, label=lst:insert_review]
-- Reviews only from verified purchases (Delivered orders)

-- Review 1: Customer 1 reviewed Product 1 (iPhone 15)
INSERT INTO Review (customer_id, target_type, target_id, 
                    rating, comment, image_url)
VALUES
    (1, 'Product', 1, 5,
     'Excellent phone! Fast performance and great camera',
     'review_iphone.jpg'),
    
    -- Review 2: Customer 1 reviewed Shop 4
    (1, 'Shop', 4, 5,
     'Very reliable seller, fast shipping',
     NULL),
    
    -- Review 3: Customer 3 reviewed Product 3 (T-Shirt)
    (3, 'Product', 3, 4,
     'Good quality cotton, true to size',
     'review_tshirt.jpg'),
    
    -- Review 4: Customer 3 reviewed Shop 5
    (3, 'Shop', 5, 4,
     'Decent shop, packaging could be better',
     NULL),
    
    -- Review 5: Customer 1 reviewed Product 5 (Headset)
    (1, 'Product', 5, 5,
     'Amazing sound quality, very comfortable',
     'review_headset.jpg'),
    
    -- Review 6: Customer 1 reviewed Shop 6
    (1, 'Shop', 6, 5,
     'Premium products, excellent customer service',
     NULL);
\end{lstlisting}